\setcounter{ex}{0}

% --- CÂU 1 ---
\begin{ex}
	Trong các công thức lượng giác sau, công thức nào \textbf{sai}?
	\choice
	{$\cos 2a = 2\cos^2 a - 1$}
	{$\cos 2a = 1 - 2\sin^2 a$}
	{\True $\cos 2a = \sin^2 a - \cos^2 a$}
	{$\sin 2a = 2\sin a \cos a$}
	\loigiai{
		Theo công thức nhân đôi, ta có $\cos 2a = \cos^2 a - \sin^2 a = 2\cos^2 a - 1 = 1 - 2\sin^2 a$.\\
		Do đó đẳng thức $\cos 2a = \sin^2 a - \cos^2 a$ là sai.
	}
\end{ex}

% --- CÂU 2 ---
\begin{ex}
	Cho $\pi < \alpha < \dfrac{3\pi}{2}$. Mệnh đề nào sau đây \textbf{đúng}?
	\choice
	{$\sin\alpha > 0$}
	{$\cos\alpha > 0$}
	{\True $\tan\alpha > 0$}
	{$\cot(\alpha + \pi) < 0$}
	\loigiai{
		Vì $\pi < \alpha < \dfrac{3\pi}{2}$ nên điểm biểu diễn của góc lượng giác $\alpha$ nằm ở góc phần tư thứ III.\\
		Khi đó $\sin\alpha < 0$ và $\cos\alpha < 0$, suy ra $\tan\alpha = \dfrac{\sin\alpha}{\cos\alpha} > 0$ và $\cot(\alpha + \pi) = \cot\alpha > 0$.
	}
\end{ex}

% --- CÂU 3 ---
\begin{ex}
	Một hình chóp ngũ giác có số đỉnh và số cạnh lần lượt là
	\choice
	{$5$ đỉnh, $10$ cạnh}
	{\True $6$ đỉnh, $10$ cạnh}
	{$6$ đỉnh, $12$ cạnh}
	{$5$ đỉnh, $8$ cạnh}
	\loigiai{
		Hình chóp ngũ giác có đáy là một ngũ giác ($5$ đỉnh, $5$ cạnh đáy) và $1$ đỉnh chóp nối với $5$ đỉnh đáy ($5$ cạnh bên).\\
		Vậy hình chóp ngũ giác có $6$ đỉnh và $10$ cạnh.
	}
\end{ex}

% --- CÂU 4 ---
\begin{ex}
	Phương trình $\sin x = 0$ có một nghiệm thuộc khoảng $\left(\dfrac{\pi}{2}; \dfrac{3\pi}{2}\right)$ là
	\choice
	{$x = \dfrac{3\pi}{4}$}
	{\True $x = \pi$}
	{$x = \dfrac{2\pi}{3}$}
	{$x = \dfrac{5\pi}{6}$}
	\loigiai{
		Ta có $\sin x = 0 \iff x = k\pi\ (k \in \mathbb{Z})$.\\
		Vì $x \in \left(\dfrac{\pi}{2}; \dfrac{3\pi}{2}\right)$ nên chọn $k = 1$, ta được nghiệm $x = \pi$.
	}
\end{ex}

% --- CÂU 5 ---
\begin{ex}
	Cho các góc lượng giác $(Ou, Ov)$, với $\text{sđ}(Ou, Ov) = -\dfrac{\pi}{3} + k2\pi, k \in \mathbb{Z}$. Góc lượng giác nào có số đo dưới đây có cùng tia đầu $Ou$ và tia cuối $Ov$?
	\choice
	{$\dfrac{\pi}{3}$}
	{$\dfrac{2\pi}{3}$}
	{\True $\dfrac{5\pi}{3}$}
	{$\dfrac{4\pi}{3}$}
	\loigiai{
		Với $k = 1 \in \mathbb{Z}$, ta có:
		$$\text{sđ}(Ou, Ov) = -\dfrac{\pi}{3} + 2\pi = \dfrac{5\pi}{3}.$$
		Do đó góc lượng giác có số đo $\dfrac{5\pi}{3}$ có cùng tia đầu $Ou$ và tia cuối $Ov$.
	}
\end{ex}

% --- CÂU 6 ---
\begin{ex}
	Cho hình chóp $S.ABCD$ có đáy $ABCD$ là hình thang ($AB \parallel CD, AB > CD$), gọi $I = AD \cap BC$. Giao tuyến của hai mặt phẳng $(SAD)$ và $(SBC)$ là đường thẳng
	\choice
	{$SA$}
	{$SB$}
	{\True $SI$}
	{$SC$}
	\loigiai{
		Ta có $S \in (SAD) \cap (SBC)$.\\
		Mặt khác $\heva{I \in AD \subset (SAD) \\ I \in BC \subset (SBC)} \implies I \in (SAD) \cap (SBC)$.\\
		Vậy giao tuyến của hai mặt phẳng $(SAD)$ và $(SBC)$ là đường thẳng $SI$.
	}
\end{ex}

% --- CÂU 7 ---
\begin{ex}
	Cho dãy số $(u_n)$, biết $\heva{u_1 &= 1 \\ u_{n+1} &= 3u_n + 1}$ với $n \ge 1$. Ba số hạng đầu tiên của dãy số đó lần lượt là những số nào dưới đây?
	\choice
	{\True $1; 4; 13$}
	{$1; 3; 10$}
	{$1; 4; 12$}
	{$1; 5; 16$}
	\loigiai{
		Từ hệ thức truy hồi của dãy số, ta tính được:
		\begin{itemize}
			\item $u_1 = 1$.
			\item $u_2 = 3u_1 + 1 = 3 \cdot 1 + 1 = 4$.
			\item $u_3 = 3u_2 + 1 = 3 \cdot 4 + 1 = 13$.
		\end{itemize}
		Vậy ba số hạng đầu tiên của dãy số là $1; 4; 13$.
	}
\end{ex}

% --- CÂU 8 ---
\begin{ex}
	Trong các dãy số được cho dưới đây, dãy số nào là một cấp số cộng?
	\choice
	{$u_n = 3^n - 1$}
	{\True $u_n = -3n + 5$}
	{$u_n = \dfrac{2}{n + 2}$}
	{$u_n = n^2 - 1$}
	\loigiai{
		Xét dãy số $(u_n)$ với $u_n = -3n + 5$. Ta có:
		$$u_{n+1} - u_n = [-3(n + 1) + 5] - (-3n + 5) = -3, \quad \forall n \in \mathbb{N}^*.$$
		Vì hiệu $u_{n+1} - u_n = -3$ không đổi nên dãy số $u_n = -3n + 5$ là một cấp số cộng với công sai $d = -3$.
	}
\end{ex}

% --- CÂU 9 ---
\begin{ex}
	Tìm tất cả các giá trị thực của $x$ để ba số $x + 1;\ x + 3;\ 2x + 3$ theo thứ tự đó lập thành một cấp số nhân.
	\choice
	{$x = 3$}
	{$x = -2$}
	{\True $x = 3$ hoặc $x = -2$}
	{$x = 2$ hoặc $x = -3$}
	\loigiai{
		Ba số $x + 1;\ x + 3;\ 2x + 3$ theo thứ tự đó lập thành một cấp số nhân khi và chỉ khi:
		$$(x + 3)^2 = (x + 1)(2x + 3) \iff x^2 + 6x + 9 = 2x^2 + 5x + 3 \iff x^2 - x - 6 = 0 \iff \hoac{x &= 3 \\ x &= -2.}$$
		Thử lại:
		\begin{itemize}
			\item Với $x = 3$, ba số là $4; 6; 9$ (lập thành cấp số nhân có công bội $q = \dfrac{3}{2}$).
			\item Với $x = -2$, ba số là $-1; 1; -1$ (lập thành cấp số nhân có công bội $q = -1$).
		\end{itemize}
		Vậy $x = 3$ hoặc $x = -2$.
	}
\end{ex}

% --- CÂU 10 ---
\begin{ex}
	Gọi $M$ và $m$ lần lượt là giá trị lớn nhất và giá trị nhỏ nhất của hàm số $y = 5\cos^2 x - 2$. Giá trị của biểu thức $P = M - m$ bằng
	\choice
	{$3$}
	{$1$}
	{\True $5$}
	{$7$}
	\loigiai{
		Với mọi $x \in \mathbb{R}$, ta có $0 \le \cos^2 x \le 1 \implies -2 \le 5\cos^2 x - 2 \le 3$.\\
		Suy ra giá trị lớn nhất $M = 3$ (khi $\cos^2 x = 1$) và giá trị nhỏ nhất $m = -2$ (khi $\cos^2 x = 0$).\\
		Vậy $P = M - m = 3 - (-2) = 5$.
	}
\end{ex}

% --- CÂU 11 ---
\begin{ex}
	Biết rằng $\sin\dfrac{\pi}{12} = \dfrac{\sqrt{a} - \sqrt{b}}{4}$ với $a, b$ là các số nguyên dương. Giá trị của biểu thức $P = a^2 - b^2$ bằng
	\choice
	{$16$}
	{\True $32$}
	{$40$}
	{$34$}
	\loigiai{
		Ta có:
		$$\sin\dfrac{\pi}{12} = \sin\left(\dfrac{\pi}{3} - \dfrac{\pi}{4}\right) = \sin\dfrac{\pi}{3}\cos\dfrac{\pi}{4} - \cos\dfrac{\pi}{3}\sin\dfrac{\pi}{4} = \dfrac{\sqrt{3}}{2}\cdot\dfrac{\sqrt{2}}{2} - \dfrac{1}{2}\cdot\dfrac{\sqrt{2}}{2} = \dfrac{\sqrt{6} - \sqrt{2}}{4}.$$
		Do đó $a = 6, b = 2$, suy ra $P = a^2 - b^2 = 6^2 - 2^2 = 36 - 4 = 32$.
	}
\end{ex}

% --- CÂU 12 ---
\begin{ex}
	Cho hình tứ diện $SABC$. Gọi $M$ là trung điểm của cạnh $SA$, $N$ là điểm trên cạnh $SB$ sao cho $SN = 3NB$ như hình vẽ sau:
	\begin{center}
		\begin{tikzpicture}[scale=0.95, font=\footnotesize, line join=round, line cap=round]
			\coordinate (A) at (0,0);
			\coordinate (C) at (3.2,0);
			\coordinate (B) at (1.3,-1.4);
			\coordinate (S) at (1.3,2.8);
			\coordinate (M) at ($(S)!0.5!(A)$);
			\coordinate (N) at ($(S)!0.75!(B)$);
			\coordinate (I) at (intersection of M--N and A--B);
			\draw[thick] (S)--(A) (S)--(B) (S)--(C) (A)--(B)--(C) (M)--(I) (B)--(I);
			\draw[dashed] (A)--(C);
			\foreach \p/\pos in {S/above, A/left, B/below left, C/right, M/above left, N/right, I/below right}
				\node[\pos] at (\p) {$\p$};
			\foreach \p in {S,A,B,C,M,N,I}
				\fill (\p) circle (1.2pt);
		\end{tikzpicture}
	\end{center}
	Trong mặt phẳng $(SAB)$, đường thẳng $MN$ cắt đường thẳng $AB$ tại điểm $I$. Tỉ số $\dfrac{IB}{IA}$ bằng
	\choice
	{\True $\dfrac{1}{3}$}
	{$\dfrac{1}{2}$}
	{$\dfrac{2}{3}$}
	{$\dfrac{1}{4}$}
	\loigiai{
		Trong mặt phẳng $(SAB)$, kẻ đường thẳng qua $B$ song song với $SA$ cắt đường thẳng $MN$ tại điểm $E$.\\
		Theo định lí Thales, ta có:
		$$\dfrac{BE}{MS} = \dfrac{NB}{NS} = \dfrac{1}{3} \quad \text{và} \quad \dfrac{IB}{IA} = \dfrac{BE}{MA}.$$
		Vì $M$ là trung điểm của $SA$ nên $MA = MS$, suy ra:
		$$\dfrac{IB}{IA} = \dfrac{BE}{MS} = \dfrac{1}{3}.$$
	}
\end{ex}
